\section{Strategic Portfolios and Customer Suitability}
\subsection{A convex allocation problem for each horizon}
For $N$ training paths and a target $w$ restored every year, the largest cumulative log loss over any contiguous interval is
\begin{equation}
Y_s(w)=\max_{0\leq u\leq v\leq H}\left\{-\sum_{t=u+1}^{v}\log(1+\rho_{s,t}^{\mathsf T}w)\right\}.
\end{equation}
The empty interval contributes zero. On this cost-free, annually restored model an ordinary maximum drawdown equals $1-\exp(-Y_s)$, but the expectation does not commute with that transform: $\ES(Y)$ does not generally give $\ES(D)$. We therefore constrain the quantity we can compute and report the resulting drawdown afterwards.

Using the variational representation of expected shortfall \cite{cvar},
\begin{equation}
\ES_\alpha(Y)=\min_z\left[z+\frac{1}{N(1-\alpha)}\sum_{s=1}^{N}(Y_s-z)_+\right],\quad \alpha=0.90,
\end{equation}
we choose the strategic anchor $w^*$ by solving
\begin{align}
\min_w\quad &-\frac{1}{NH}\sum_{s=1}^{N}\sum_{t=1}^{H}\log(1+\rho_{s,t}^{\mathsf T}w)\\
\text{subject to}\quad &\ES_{0.90}(Y(w))\leq qR_{\rm ref},\quad w\geq0,\quad\mathbf{1}^{\mathsf T}w=1.
\label{eq:strategy}
\end{align}

For each horizon, $R_{\rm ref}$ is the training log-drawdown expected shortfall of a stated reference allocation: $(.20,.30,.50)$ for Conservative, $(.50,.35,.15)$ for Balanced and $(.75,.20,.05)$ for Growth. The multiplier $q\in\{.75,.85,.95\}$ scales the training budget and is chosen in Section~3.2, while the validation target stays at $R_{\rm ref}$. Conservative additionally requires $w_{\rm bill}\geq0.40$. The three products thus differ both in risk budget and, for Conservative, in a liquidity floor.

Negative log growth and the maximum of interval log losses are convex on the valid-return domain, so \eqref{eq:strategy} is a convex subproblem. We solve it with a Kelley cutting-plane linear program: supporting hyperplanes generate objective lower bounds, feasible weights upper bounds, and the loop stops at a gap of $10^{-6}$, with risk tolerance $10^{-8}$ and at most 300 iterations, starting from a 231-point 5\% simplex grid that also supplies a comparison objective. Primal solutions and dual multipliers are saved so that feasibility and duality balance can be re-checked. An infeasible risk budget is recorded and never repaired by quietly relaxing it. These certificates cover the strategic subproblem only; they say nothing about the optimality of the trading and escalation layer.

\subsection{Development selection before freezing}
Selection happens before the holdout is touched. Nested validation uses five outer folds of complete paths, 96 scenarios to train and 24 to validate, and three inner folds of 64 and 32 inside each outer training set, so every choice is scored on data the candidate was not fitted to. Candidates pair $q\in\{.75,.85,.95\}$ with execution curvature $\gamma\in\{4,12\}$. Comparing them on held-out outcomes rather than on training fit limits the selection bias that repeated comparisons introduce \cite{cawley}.

A candidate is feasible on validation only if four conditions hold. Pooled log-drawdown expected shortfall at 90\% must not exceed $R_{\rm ref}$ plus 0.005 in log units; that slack is fixed in advance and does not follow $q$. The median and the 10th percentile of annual real growth must not fall more than 0.001 and 0.002 below fee-paying annual rebalancing to the same anchor. Reviews, and hence the trades they permit, must be requested in at most 10\% of validation scenarios. Among feasible candidates we maximize the mean of $\log(1+G)$. If none is feasible, the least-violating candidate is retained as a diagnostic and marked unusable rather than presented as a selection. All three products selected $q=0.95$ and $\gamma=4$ on the full development set, after which the anchors and the customer assignments were frozen.

\subsection{Frozen allocations and an explicit refusal option}
Table~\ref{tab:weights} gives the frozen anchors at the five horizons that appear in the questionnaire. A product name does not imply one allocation for every customer: Conservative is nearly invariant to the horizon, holding 70--74\% in bills throughout, while Balanced and Growth move bills into bonds as the horizon lengthens, and the 30-year Growth solution reaches the boundary with no bills at all. The response is not monotone in the middle of the range, because each anchor is fitted separately at its own horizon. These are model outputs under the stated data and constraints, not general investment prescriptions.

\begin{table}[h!]
\centering
\caption{Frozen strategic equity / bond / bill weights (\%) from the full-development refit; rounding may make totals differ from 100.}
\label{tab:weights}
\begin{tabular}{rlll}\toprule
$H$ & Conservative & Balanced & Growth\\\midrule
\tableinput{generated/weights.tex}
\bottomrule\end{tabular}
\end{table}

For horizon $H$ and tolerance $L$ we count out-of-fold paths whose complete execution policy breaches $D>L$; with $k$ breaches among $n=120$, the one-sided Clopper--Pearson upper bound is
\begin{equation}
U=\operatorname{Beta}^{-1}(1-\delta;k+1,n-k),\qquad
\delta=0.05/440,
\end{equation}
with $U=1$ when $k=n$. The family of 440 is fixed before screening and covers five profiles, 11 distinct product-or-programme choices, four liquidity conventions and two cost levels. Automatic service requires $U\leq0.20$ at the base cost and at doubled cost alike. Because the out-of-fold policies are refitted on overlapping training sets, $U$ is a conservative screening diagnostic rather than a coverage guarantee for the final refitted policy, and the 20\% ceiling is a service assumption, not a stated customer preference.

Two further rules follow from the questionnaire. A high-liquidity profile requires a strategic bill weight of at least 40\% in every relevant development fold; the condition is checked for any eligible product, so such a profile is not forced into the product named Conservative, and the target is not a cash floor maintained after market drift. Growth is offered only from a 15-year horizon upwards, and P90 drawdown must be ordered Conservative $\leq$ Balanced $\leq$ Growth. Among the pairs that clear every rule, the one with the highest out-of-fold median growth is recommended; if none clears them, the profile goes to human review.

Table~\ref{tab:suitability} applies these rules pair by pair. For A and D no product comes close, so referral is the model's answer rather than a default. B and C clear only through Conservative, and C's Balanced alternative fails the ceiling on 17 breaches. E has two qualifying pairs; Balanced wins on out-of-fold median growth, 4.22\% against 2.71\%. All three served profiles also clear the doubled-cost test.

\begin{table}[h!]
\centering
\caption{Screening of every profile--product pair on development data: breaching paths $k/120$ and the bound $U$ (\%). Bold marks a pair that also clears the liquidity, horizon, ordering and doubled-cost rules.}
\label{tab:suitability}
\begin{tabular}{lcclccccccc}\toprule
Profile & $H$ & Tol.\,(\%) & Liquidity & \multicolumn{2}{c}{Conservative} & \multicolumn{2}{c}{Balanced} & \multicolumn{2}{c}{Growth} & Assigned\\
\cmidrule(lr){5-6}\cmidrule(lr){7-8}\cmidrule(lr){9-10}
 & & & & $k/120$ & $U$ & $k/120$ & $U$ & $k/120$ & $U$ & \\\midrule
\tableinput{generated/suitability.tex}
\bottomrule\end{tabular}
\end{table}
